\documentclass[11pt,a4paper]{article}
\usepackage[utf8]{inputenc}
\usepackage[german,english]{babel}
\usepackage[margin=0.8in]{geometry}
\usepackage{xcolor}
\usepackage{tikz}
\usetikzlibrary{shapes.geometric, shapes.symbols, positioning, arrows.meta, shadows, calc}
\usepackage[most]{tcolorbox}
\usepackage{enumitem}
\usepackage{hyperref}
\usepackage{booktabs}
\usepackage{array}
\usepackage{titlesec}

% 🎨 Custom Colors (FinTech Premium Palette)
\definecolor{primarynavy}{HTML}{0B0F19}     % Deep dark blue
\definecolor{accentblue}{HTML}{2563EB}      % Royal Blue
\definecolor{accentgreen}{HTML}{10B981}     % Emerald Green
\definecolor{accentorange}{HTML}{F97316}    % Coral Orange
\definecolor{accentpurple}{HTML}{A855F7}    % Amethyst Purple
\definecolor{accentcyan}{HTML}{06B6D4}      % Cyber Cyan
\definecolor{accentred}{HTML}{EF4444}       % Crimson Red
\definecolor{bglight}{HTML}{F8FAFC}        % Off-white background
\definecolor{bordergray}{HTML}{E2E8F0}     % Slate border

\titleformat{\section}
  {\color{primarynavy}\normalfont\Large\bfseries}
  {\thesection}{1em}{}[{\color{accentblue}\titlerule[2pt]}]
\titleformat{\subsection}
  {\color{accentblue}\normalfont\large\bfseries}
  {\thesubsection}{1em}{}

\setlist[itemize]{label=\color{accentblue}$\bullet$}
\setlist[enumerate]{label=\color{accentblue}\arabic*.}
\setlength{\parindent}{0pt}
\setlength{\parskip}{8pt}

\hypersetup{
    colorlinks=true,
    linkcolor=accentblue,
    urlcolor=accentblue
}

\begin{document}
\selectlanguage{english}

% ==========================================
% 👑 HEADER BANNER
% ==========================================
\begin{tcolorbox}[
    colback=primarynavy,
    colframe=accentblue,
    arc=8pt,
    boxrule=2.5pt,
    left=20pt,
    right=20pt,
    top=25pt,
    bottom=25pt,
    coltext=white
]
    \centering
    {\huge\bfseries Kubernetes (K8s) \& Microsoft Azure \par}
    \vspace{8pt}
    {\Large\bfseries Cloud Orchestration \& Backend Safeguards Guide \par}
    \vspace{10pt}
    \color{accentcyan}\hrule height 1.5pt
    \vspace{8pt}
    {\small Student Project Reference -- \textbf{Andy Rugero Ndayambaje} \par}
\end{tcolorbox}

% ==========================================
% 🌐 CORE SYNERGY
% ==========================================
\section{The Core Cloud Synergy: Azure \& Kubernetes}

In a modern production application like \textbf{Vindobona Pro FinTech}, the physical servers (cloud) and the deployment manager (orchestrator) work hand in hand.

\begin{tcolorbox}[
    colback=bglight,
    colframe=accentorange,
    arc=6pt,
    boxrule=1.5pt,
    title=\textbf{\color{white}💡 The Real-World Analogy}
]
    Imagine building a modular residential complex:
    \begin{itemize}
        \item \textbf{Azure is the Landlord \& Real Estate Developer:} They supply the plot of land, water pipes, electrical grid, and building foundations (Virtual Machines, physical network, raw databases).
        \item \textbf{Kubernetes is the Construction Foreman:} They do not own the land, but they coordinate everything on it. They decide how many modular apartments to place, arrange them for optimal sunlight, and immediately replace a module if it gets damaged.
        \item \textbf{Docker Containers are the Prefabricated Apartments:} Independent, identical, self-contained units (Express API, React Frontend) ready to drop onto the foundation.
    \end{itemize}
\end{tcolorbox}

\subsection{Distinction of Responsibilities}
\begin{table}[h]
\centering
\small
\begin{tabular}{lp{6.5cm}p{6.5cm}}
\toprule
\textbf{Feature} & \textbf{Microsoft Azure (Hardware Layer)} & \textbf{Kubernetes (Orchestration Layer)} \\
\midrule
\textbf{Role} & Provides hardware \& cloud services. & Manages, monitors, and scales containers. \\
\textbf{Primary Asset} & Virtual Machines (VMs), storage, physical network nodes. & Software controller running inside virtual clusters. \\
\textbf{Scaling} & Provisions new physical virtual computers. & Launches new app replicas on existing computers. \\
\textbf{Health Action} & Alerts if a physical hard drive/node dies. & Restarts app containers inside VMs when they crash. \\
\bottomrule
\end{tabular}
\caption{Azure vs. Kubernetes Division of Work}
\end{table}

\newpage

% ==========================================
% 📊 ARCHITECTURE DIAGRAM
% ==========================================
\section{Deployment Architecture Diagram}

Below is the visual schematic showing the request lifecycle from the browser to the backend server replicas managed by Kubernetes on Microsoft Azure:

\vspace{0.8cm}

\begin{figure}[h]
\centering
\begin{tikzpicture}[
    node distance=1.8cm and 1.8cm,
    clientNode/.style={circle, draw=accentblue!80, fill=accentblue!10, thick, minimum size=2.2cm, align=center, font=\scriptsize\bfseries, drop shadow},
    lbNode/.style={rectangle, draw=accentorange!80, fill=accentorange!15, thick, minimum width=2.8cm, minimum height=1.2cm, align=center, rounded corners=4pt, drop shadow},
    podNode/.style={rectangle, draw=accentpurple!80, fill=accentpurple!10, thick, minimum width=2.4cm, minimum height=1.3cm, align=center, rounded corners=4pt, drop shadow},
    dbNode/.style={cylinder, draw=accentgreen!80, fill=accentgreen!10, shape border rotate=90, minimum width=2.2cm, minimum height=1.4cm, align=center, aspect=0.25, drop shadow},
    arrow/.style={-Stealth, ultra thick, draw=primarynavy!70}
]

    % Nodes
    \node (client) [clientNode] {Client Browser\\(Vindobona UI)};
    \node (lb) [lbNode, right=of client] {\textbf{K8s Service}\\(Load Balancer)};
    
    \node (pod2) [podNode, right=of lb] {Express API\\(Replica 2)};
    \node (pod1) [podNode, above=0.6cm of pod2] {Express API\\(Replica 1)};
    \node (pod3) [podNode, below=0.6cm of pod2] {Express API\\(Replica 3)};
    
    \node (db) [dbNode, right=of pod2] {PostgreSQL\\(Azure DB)};

    % Boundary Boxes
    \draw[dashed, draw=accentcyan, thick] ($(pod1.north west)+(-0.3,0.4)$) rectangle ($(pod3.south east)+(0.3,-0.4)$) node[pos=0.08, above, font=\tiny\bfseries\color{accentcyan}] {Kubernetes Cluster (AKS)};
    \draw[dashed, draw=accentblue!50, thick] ($(lb.north west)+(-0.2,1.8)$) rectangle ($(db.south east)+(0.4,-1.8)$) node[pos=0.96, below, font=\tiny\bfseries\color{accentblue}] {Microsoft Azure Cloud Zone};

    % Connections
    \draw[arrow] (client) -- node[above, font=\tiny\bfseries] {HTTPS :443} (lb);
    \draw[arrow] (lb.east) -- ++(0.5,0) |- (pod1.west);
    \draw[arrow] (lb.east) -- (pod2.west);
    \draw[arrow] (lb.east) -- ++(0.5,0) |- (pod3.west);
    
    \draw[arrow] (pod1.east) -- ++(0.5,0) |- (db.north);
    \draw[arrow] (pod2.east) -- (db.west);
    \draw[arrow] (pod3.east) -- ++(0.5,0) |- (db.south);

\end{tikzpicture}
\caption{Load-balanced deployment schematic of Vindobona Pro FinTech on Azure}
\end{figure}

\vspace{0.5cm}

% ==========================================
% 🛡️ BACKEND SAFEGUARDS
% ==========================================
\section{How Kubernetes Protects the Backend Express API}

Running a backend server without orchestrators is highly risky. Kubernetes acts as an automated operations team, protecting the server through four specific mechanisms:

\subsection{1. Self-Healing (Automated Container Recovery)}
If the Express API crashes due to an unhandled promise rejection, memory leaks, or execution failures:
\begin{itemize}
    \item \textbf{Without K8s:} The application remains offline until a developer SSHs into the machine and restarts the process manually.
    \item \textbf{With K8s:} Kubernetes monitors the container container state. If the process terminates, it destroys the failed container and spawns a clean, new container replica **in less than 3 seconds**.
\end{itemize}
\textbf{Real-World Example:} A user sends malformed transaction payload causing Node.js to throw a fatal error. The server crashes; Kubernetes intercepts it immediately and boots a new instance, preventing downtime for other users.

\subsection{2. Load Balancing (Traffic Equalization)}
During high-traffic situations (e.g., end-of-month paydays):
\begin{itemize}
    \item \textbf{Without K8s:} All incoming API requests land on a single Node.js thread, causing response lag, high CPU load, and eventual server timeout.
    \item \textbf{With K8s:} Kubernetes routes traffic dynamically across the **3 active replicas**. If Replica 1 is busy executing a calculation, the next request is automatically routed to Replica 2 or 3.
\end{itemize}

\newpage

\subsection{3. Zero-Downtime Updates (Rolling Deployments)}
When code changes (like adding a feature to \texttt{server.js}) are deployed:
\begin{itemize}
    \item \textbf{Without K8s:} The backend process must be stopped, updated, and restarted. During these minutes, users face a \textit{503 Service Unavailable} error page.
    \item \textbf{With K8s:} Kubernetes conducts a **Rolling Update**:
    \begin{enumerate}
        \item It launches the new version as a separate container.
        \item Once the new container passes health checks, Kubernetes routes traffic to it.
        \item It destroys the old container replica.
        \item This sequence repeats one-by-one until all 3 replicas are updated without dropping a single connection.
    \end{enumerate}
\end{itemize}

\subsection{4. Intelligent Liveness \& Readiness Probes}
Using the Express health endpoint (\texttt{/api/health}):
\begin{itemize}
    \item \textbf{Without K8s:} If the Express server hangs in an infinite loop or loses its database pool connection, it remains running but completely unresponsive.
    \item \textbf{With K8s:} Kubernetes pings the \texttt{/api/health} endpoint every 10 seconds. If the API returns a status other than \texttt{200 OK} (or fails to respond), Kubernetes labels the container as unhealthy, removes it from the load balancer, and deploys a healthy instance.
\end{itemize}

\begin{tcolorbox}[
    colback=bglight,
    colframe=accentgreen,
    arc=6pt,
    boxrule=1.5pt,
    title=\textbf{\color{white}🔒 Summary for Your Presentation}
]
    \centering
    \small\itshape
    "Microsoft Azure serves as our hardware foundation and scalable cloud provider. Kubernetes acts as our deployment manager on top, ensuring that our Express backend remains resilient through automated self-healing, balanced workloads across 3 replicas, zero-downtime rolling updates, and active health monitoring."
\end{tcolorbox}

\end{document}
